\section{把慢变量放回年序：压力如何一再转移}
\label{sec:synthesis-chronology-pressures}

综合章不能把三百多年历史从年序中抽走，写成几条永恒主题。迁徙、家户、制度、边境、宗教和文本之所以重要，正在于它们在不同转折点承受不同压力。西晋统一时需要接收吴地，八王之乱时需要维持都城与宗室兵权，永嘉之后需要安置与登记流动者，东晋北伐时需要把江南供给送往淮河以北，北魏统一时需要把征服转为户籍和田役，迁都与六镇危机时需要重新平衡中心与边地，南北末期则要在战争、寺院、城市和外交之间筹集资源。时间顺序使这些问题不至于互相替代。

\subsection{二八〇年的统一：接收的压力}

孙吴降服使西晋在名义上结束三国并立。这个结果可作为全国统一的起点，却不能把它当成国家能力已经完成的证明。建业以东的州郡、旧官员、江面船运和地方社会需要被接收；北方中央则要决定怎样把新得地区纳入赋役、任命和裁判。统一的好处是减少了长期边界战争，成本是行政与财政必须覆盖更大的差异。若接收不能稳定，后来的内乱便会更容易把地方资源从朝廷手中抽走。

八王之乱将这种脆弱暴露在洛阳。宗室拥有军事资源，宫廷掌握诏令与信息，地方军队又要依赖粮饷和道路；这些条件在连续清洗中相互撕裂。不能用“制度腐败”一句解释所有行动，也不能把各王的选择全写成私人野心。更具体的问题是，谁能让军队进入京师、谁能供给它、谁能以皇帝名义发布命令、谁又能在城市之外组织支持。都城失去协调能力后，西晋的统一遗产不再能约束多方行动。

\subsection{三一一年之后：安置的压力}

洛阳失守和北方政局瓦解使人、田、文书与军队进入持续移动。东晋在江南建立并非从空白处创造新社会，而是要让旧籍、侨置、地方强宗和水路供给在新的皇统下共存。南渡的压力首先是安置：迁来者怎样取得住处与职位，本地居民怎样面对新的征发与竞争，朝廷怎样在没有完整旧档的条件下建立可用的登记。

这种压力随北伐而改变。祖逖、桓温、刘裕等人的军事行动会使收复、投附与再次逃散成为现实可能，前线一旦延伸，后方的江淮、荆州和建康就要承担粮运。北伐成功的象征意义很强，实际接收却需要更久。读者应把“未能北伐成功”从道德性评价转回资源问题：道路、船运、地方接应、军队纪律与朝廷内部协调是否足以让一时的突破变成长久控制。

\subsection{四三九年前后：统一北方的压力}

北魏先后取中山、统万、龙城、姑臧，将北方主要割据政权逐步纳入更大框架。四三九年灭北凉可作为北方军事整合的标志，却不应写成河西、代北、河北与关中从此具有同样的治理密度。每一新得地区都带来旧官、降众、城镇、道路与宗教网络；军队能够到达的地方，未必已能稳定征税、裁判或登记。

均田、三长、班禄等改革正是这一压力的回答。它们试图把征服的结果转成能持续的田、户、役和基层责任，而不是继续依赖临时军队。太武帝禁佛也应放在国家能力的紧张处理解：政策触及寺院、僧尼与资源，却不能被还原为单纯财政或单纯信仰。北魏此时面对的，是怎样让多个地区与多种中介都可被国家看见；政策的意图清楚，实际执行则仍要分地点和时段判断。

\subsection{四九四年至五二八年：中心转移的压力}

迁都洛阳改变了北魏的政治地理。它使朝廷更接近中原农业区、旧都礼制和新的城市营造，也使原来依赖平城与北方军镇的力量重新估量自己的位置。改姓、服饰、礼仪和婚姻政策在公共层面塑造新的规范，但并不让所有地方关系同时重写。迁都既是王权追求更有效统治的选择，也把供给、升迁与文化声望的机会向新中心集中。

六镇起事表明边地压力不会因都城更华丽而消失。镇户的生活、军粮、地位与将领关系，在长期累积后成为动乱条件；起事之后的军人、流民与地方势力重新流入河北、河东、关中和城市网络。北魏的裂解并非一场宫廷斗争的简单外溢，而是中心与边地、登记与供养、军队与家户原本未被解决的矛盾在特定时刻失去承受力。

\begin{causalchain}[从迁都到分裂：不要把先后当作单一因果]
洛阳迁都改变了资源、仪礼与官职的重心；边镇的供给、地位与地方关系则有更长的积累。六镇动乱发生在这些条件相互作用之中，而不是由某一服饰、姓氏或单一诏令自动引发。分裂之后，东魏北齐与西魏北周分别吸纳不同区域的军人、城镇与财政资源，才使原来的危机成为新的政权格局。
\end{causalchain}

\subsection{五三四年至五七七年：竞争的压力}

东西魏分立后，河北与关中的资源差异进入长期战争。东魏北齐拥有邺及河北的城市、人口和手工业，西魏北周则以关中核心、军府与灵活联盟求得生存。两方都不能只靠一场战役决定命运：军队要持续获得粮食，官员要使征收与裁判继续，寺院、地方豪强和军人家户也都在资源分配中拥有位置。制度得失不在于某种名目更古雅，而在于它能否让有限资源被重复调动而不立刻毁坏社会关系。

北周灭北齐改变了这场竞争，却把接收的压力带回前台。河北旧有军政和城市网络必须被利用、安抚或重组；宗教政策、官制调整和对外关系也都需要物资和人力。胜者若只留下新的官号而无法让地方生活恢复可预期，统一就会停留在军报上。北周的短暂、隋的继起，说明制度与资源网络比一个统治者的寿命更能决定接收能否延续。

\subsection{五八九年：统一的压力仍在继续}

陈亡使南北长期并立的最高政权结构结束。隋军南下所依赖的北方军政、江淮道路与接收安排，是此前数十年竞争积累的结果；陈朝州郡、官员、仓储、水路与寺院则成为新政权必须处理的遗产。对史书而言，受降可以标示一个终点；对家户、地方官和城市而言，真正的变化还要在登记、调运、婚姻、市场和裁判中展开。

因此，本卷的结尾不应把隋写成解决一切的答案。再统一确实减少了南北对峙的最高政治边界，也提供了重整州郡与交通的机会；它同时把更大范围的人口、宗教资源、边军与地方差异放入同一政府的责任之下。两晋南北朝留下的不是一段“分裂终于被克服”的空白，而是一整套如何接收流动者、维持家户、利用道路、处理寺院并让制度经过地方中介的难题。

\begin{debate}[能否用“乱世”概括二八〇至五八九年]
“乱世”能提示战争与政权更替的密度，却容易遮住西晋统一、东晋与南朝的长期行政、北魏的北方整合、城市营造、文本生产和地方生活的连续性。更准确的说法是：这是多次统一、崩解、接收与再组织交错的时期。它的秩序并非缺席，而是在不同区域、不同制度和不同人群中以不相同的速度形成与破裂。
\end{debate}

将压力放回年序，能避免两种相反的写法：一是把每次政变都看成孤立的人物戏，二是让迁徙、宗教或制度漂浮成与事件无关的概念。每一个转折都改变了谁承担粮、役、路、城与记忆的成本；每一项慢变量又使某些政治选择成为可能或不可能。两者共同构成这个时代真正连续的编年。
